% Part I — अध्याय १: संत दरिया बाबा के जीवन यात्रा
\section*{अध्याय १: संत दरिया बाबा के जीवन यात्रा}
\phantomsection\addcontentsline{toc}{section}{अध्याय १: संत दरिया बाबा के जीवन यात्रा}

\begin{quote}
जइसे पेड़ आ नदिया के पानी\\
दूसर के भलाई खातिर होला,\\
संत के स्वभाव भी दूसर के दुख हर के\\
आनंद देल बा।\footnote{ज्ञान रतन, चौपाई १६०५, द.ग्रं., पृ. २२३}
\end{quote}

\subsection*{१.१ जन्म आ प्रारंभिक जीवन}
\phantomsection\addcontentsline{toc}{subsection}{१.१ जन्म आ प्रारंभिक जीवन}

दु सौ बरिस से कुछे बेसी पहिले, दरिया बाबा — निस्संदेह सबसे उच्च कोटि के संत — भारत में रहत रहलें आ फिर भी हमनी ओहन के बारे में बहुत कम जानित बानी। एह में कवनो अनहोनी नइखे, काहे कि संसार आमतौर पर संतन के न सिरिफ उदासीनता आ उपेक्षा से देखे ला, बल्कि अक्सर तिरस्कार आ बैरभाव से देखे ला। सौभाग्य से, दरिया के कुछ अंतरंग शिष्य रहलें जिनकर संगे ऊ कबो-कबो आपन जीवन के कुछ घटना साझा करत रहलें, आपन उपदेश समझावत घरी। कम से कम एगो शिष्य अक्सर दरिया बाबा के बात सभ के लिख लेत रहलें। दरिया बाबा खुद शायद बहुत कम लिखलें। कुछ पद के छोड़ के जे उनकर तत्कालीन शिष्य लोग मरणोपरांत जोड़लक, उनकर लगभग सगरी रचना उनकर उक्ति सभ के लिखल संकलन हवे। उनकर जीवन के निम्नलिखित विवरण मुख्यतः दरिया बाबा के एह कृति सभ पर आधारित बा, जे में से अधिकतर अब भी हस्तलिखित पांडुलिपि के रूप में बाड़ी।

कुछ पांडुलिपि के अंत में साफ़ लिखल बा कि दरिया बाबा विक्रम संवत १८३७ (सन् १७८० ई.) में भादो बदी चतुर्थी के दिन — जे शुक्रवार रहे — आपन शरीर छोड़लें।\footnote{दरिया सागर, ज्ञान दीपक आ सहस्राणी के समापन पद।} एह तिथि पर संदेह के कवनो कारण नइखे। आम धारणा बा कि दरिया १०६ बरिस जीवित रहलें। तदनुसार, उनकर जन्म विक्रम संवत १७३१ (सन् १६७४ ई.) में भइल रहे। ई मान्यता धर्मेंद्र ब्रह्मचारी शास्त्री\footnote{ध.ब्र. शास्त्री, द.ग्रं. भाग १, पृ. ५} आ बेलवेडियर प्रेस, इलाहाबाद के सेंट सीरीज के संपादक\footnote{\textit{दरिया सागर} (इलाहाबाद: बेलवेडियर प्रेस, तीसरा संस्करण, १९७०), पृ. १} द्वारा स्वीकारल गइल बा। दरिया संप्रदाय के कुछ अनुयायी लोग विक्रम संवत १६९१ (सन् १६३४ ई.) के दरिया के जन्म तिथि माने ला। बाकिर एह से दरिया के उमिर १४६ बरिस होखी जे मुश्किल से संभव बा। दरिया अपना रचना में पंद्रहवीं, सोलहवीं आ सत्रहवीं सदी के कई संत आ रहस्यवादी लोग के उल्लेख करे लें, जइसे कबीर (१३९८-१५१८), धर्मदास (कबीर के उत्तराधिकारी), नानक (१४६९-१५३८), नामदेव (१३९८-१५१८), मीरा बाई (१४९८-१५४६) आ गोस्वामी तुलसीदास (१५३२-१६२३), आ बाद के संतन के कवनो जिक्र ना करे लें — ई बात ऊपर बतावल तिथि (१६७४-१७८०) के समर्थन करे ला। एह से पता चले ला कि ऊ एह संसार से खाली ३८ बरिस पहिले विदा भइलें, आगरा के प्रसिद्ध संत सेठ शिव दयाल सिंह जी (सोआमी जी, १८१८-१८७८) के जन्म से।

दरिया के जन्म के जगह, उनकर पिता के नाम आ पेशा के साफ़ जिक्र दरिया के रचना \textit{मूर्ति उखाड़} में मिले ला। दरिया के संदर्भ में ओह में कहल गइल बा:

\begin{quote}
पीरू दर्जी के घर में,\\
धरकंधा के निवासी,\\
ई पवित्र पुरुष (दरिया) जन्म लेले बाड़े।\footnote{मूर्ति उखाड़, पांडु., पृ. ८}
\end{quote}

धरकंधा गाँव बिहार के रोहतास जिला में पड़े ला आ आरा से लगभग ५२ मील, डुमराँव से लगभग २६ मील आ सूर्यपुरा से लगभग ६ मील दूर बा। दरिया बाबा के जन्मभूमि होखे के नाते, ई आज दरिया संप्रदाय के मुख्य केंद्र हवे जेकर लगभग १५० उपकेंद्र बिहार, उत्तर प्रदेश आ नेपाल में बाड़ी।

फ्रांसिस बुकानन, जे १८०९-१८१० के बीच शाहाबाद जिला (वर्तमान में भोजपुर आ रोहतास जिला) के भ्रमण कइलें, दरिया के निधन के मात्र ३० बरिस बाद, दरिया के एगो मुसलमान दर्जी बतवलें।\footnote{शाहाबाद रिपोर्ट, पृ. २२०} अंग्रेजी में दरिया के एकमात्र अउरी उल्लेख बी.बी. मजुमदार के छोटका लेख "द टेलर सेंट ऑफ़ बिहार" में बा, जे ११ सितंबर १९३५ के बिहार के अंग्रेजी अखबार \textit{सर्चलाइट} में छपल रहे। एह शुरुआती विवरण से ई निष्कर्ष निकले ला कि दरिया एगो मुसलमान दर्जी के परिवार में जन्मल रहलें।

बाकिर एह बात पर कुछ विवाद बा कि ई पीरू दर्जी केन रहलें आ ऊ मुसलमान वंश के रहलें कि ना। पंडित सुधाकर द्विवेदी के अनुसार, दरिया के जन्म एगो मुसलमान माता से भइल जे बादशाह औरंगजेब के बेगम के दर्जी के बेटी रहली। दरिया के पिता पूरन शाह, एगो प्रतिष्ठित क्षत्रिय, के ओहसे बियाह करे के पड़ल ताकि औरंगजेब द्वारा अपना भाई लोग के फाँसी से बचा सके।\footnote{द.ग्रं. भाग १, पृ. ८} दूसर मत के अनुसार, दरिया के जन्म उनकर पिता के पहिली पत्नी से भइल जे हिंदू रहली।\footnote{ई मत बेलवेडियर प्रेस के सेंट सीरीज के संपादक द्वारा व्यक्त कइल गइल — द.ग्रं. भाग १, पृ. ८} धर्मेंद्र ब्रह्मचारी शास्त्री बतावे लें कि दरिया संप्रदाय के साधु चतुरी दास ओह लोग के दरिया के वंशावली दिहलें जेकर अनुसार दरिया के पिता के मूल नाम पृथु देव सिंह, उपनाम पूरन शाह रहे। ई हिंदू नाम बदल के पीरू शाह भ गइल, जे इस्लाम धर्म स्वीकार करे के संकेत देला। पृथु देव सिंह के पिता के नाम सुमेर सिंह, उनकर पिता सूरत चंद्र सिंह आ उनकर पिता रणजीत नारायण सिंह बतावल गइल — ई सब स्पष्ट हिंदू नाम हवे।\footnote{द.ग्रं. भाग १, पृ. ९} धर्मेंद्र ब्रह्मचारी बतावे लें कि ऊ धरकंधा जा के दरिया बाबा के वंशज मेघबरन दास से बात कइलें। यद्यपि मेघबरन दास खाली चौथी पीढ़ी तक के पूर्वजन के नाम बता सकलें, सब नाम स्पष्ट हिंदू रहलें।\footnote{द.ग्रं. भाग १, पृ. ११} कहल जाला कि दरिया के पैतृक स्थान राजापुर रहे जे धरकंधा से लगभग १० मील दूर बा, आ धरकंधा दरिया के माता के मायका रहे। बाकिर एह बात में कवनो संदेह नइखे कि दरिया धरकंधा में जन्मलें आ पलल-बढ़लें आ अपना अधिकतर समय एहिजा बितवलें। ई सवाल कि ऊ हिंदू परिवार में जन्मलें या मुसलमान में, दरिया खातिर कवनो मायने ना रखे, काहे कि ऊ धार्मिक आ जातिगत भेद के कवनो महत्व ना देत रहलें। वास्तव में, ऊ खुलेआम घोषित कइलें कि संतन के कवनो जाति या धर्म नइखे।\footnote{जाति न पूछो, पंथ न पूछो, पूछो ग्यान बिचार। मुक्ति पदवी पा के संत, जाति-बरन से न्यार॥ (\textit{भक्ति हेतु}, सखी १६, द.ग्रं. भाग २, पृ. २८९) \\ हिंदू मुसलमान दूनो एक, जे शब्द धुन के करें टेक। समझ बूझ के कहौ पुकार, सब जीवन के प्रभु अधिकार॥ (\textit{दरिया सागर}, चौपाई ६१०-६११, द.ग्रं. भाग २, पृ. ६१)} हिंदू आ मुसलमान, बिना कवनो भेदभाव के, ओह लोग द्वारा शिष्य के रूप में स्वीकार कइल जात रहलें आ ऊ दुनो के धार्मिक ग्रंथन से भली-भाँति परिचित रहलें, जइसन कि उनकर रचना में ओह सब के प्रचुर संदर्भ से स्पष्ट बा।

दरिया के रचना सभ ओह लोग के सत्य पुरुष (परम सत्ता) के अवतार या पुत्र के रूप में चित्रित करे लीं। ओह लोग के एह संसार में सत्य पुरुष द्वारा भेजल गइल रहलें ताकि पीड़ित आत्मा सभ के मन आ माया के चंगुल से छुड़ा के ओह लोग के मूल घर — सत लोक या छप लोक — में वापस ले जा सके। दरिया दावा करे लें कि ऊ एह दुनिया में अठारह बेर आ चुकल बाड़ें, कबो लगभग अज्ञात रह के आ कबो कुछ प्रकट हो के।\footnote{अठारह जन्म हम भए, पाँच में गावल एहि (शब्द) के गुन। कबहुँक रहलौं ढाँकि-ढँपाइ, कबहुँक कहि रहलौं कछु अपन। आपन जीवन जइसन भइल, आपन हम एही लिखलौं। (\textit{ज्ञान दीपक}, पांडु., पृ. ३८)}

एहिजा प्रस्तुत विवरण मुख्यतः दरिया के रचना \textit{ज्ञान दीपक} पर आधारित बा आ एकर पूरक \textit{दरिया सागर}, \textit{भक्ति हेतु} आ \textit{मूर्ति उखाड़} जइसन दूसर रचना सभ बाड़ी।

दरिया के जन्म बड़ी धूमधाम से मनावल गइल। बच्चा के एक महीना के भइला पर एगो संत ओह जगह आइलें आ माता बच्चा के ओह लोग के सामने ले आइली। ई संत बच्चा के सिर से पैर तक ध्यान से देखलें, माता से कहलें कि ओह बच्चा के बहुत अच्छा से देखभाल करीं, आ बच्चा के नाम दरिया रखलें। ई संत कोई आरु ना रहलें बल्कि बाद में दरिया के सतगुरु के रूप में पहिचानल गइलें। दरिया कबो ओह लोग के मनुष्य के रूप में ना देखलें आ हमेशा प्रभु या सत्य पुरुष कह के संबोधित कइलें।

जब दरिया खाली नौ बरिस के रहलें, उनकर माई-बाप तत्कालीन रिवाज के अनुसार ओह लोग के बियाह कर देलन। दरिया कहतारे कि ऊ ओह अवसर के आडंबर शायद ही समझ पवलें, बाकिर ऊ टिप्पणी करे लें कि जे ओह लोग के संग बियाह में बँधली, ऊ पहिले कइल गंभीर तपस्या के फलस्वरूप एह संबंध में आइली।\footnote{बियाह करा दिहलन बिनु बार, न समझलौं हम ठगी सिगार। जेइ करम कीन्ह तप पहिलेहूँ, आइ संग लागल हमरेहूँ। एहि तरह मन बनावल माई, जेइ जगत के चलन चलाई। (\textit{ज्ञान दीपक}, पांडु., पृ. ३९७)}

दरिया के बचपन मासूमियत में बीतल; बाकिर जइसे-जइसे ऊ बड़ होखत गइलन, संसारिक मोह-माया विकसित होखे लागल, काहे कि जइसन ऊ कहतारे, उनकर भीतरी शब्द धुन अब भी सुप्त रहे।\footnote{बालकपन कछु बीति गवनू, शब्द सुनत नाहिं तब धुनू। संसार मोह लगावा कछु समुझि, मोर सब समुझि मन मानयो अप्प। (\textit{दरिया सागर}, चौपाई २४-२५, द.ग्रं. भाग २, पृ. ४)} जब ऊ सोलह बरिस के भइलन, ऊ बहुत उदास रहे लागलन। सपना में ऊ अक्सर पद रचना करत रहलें आ जागे पर ओह लोग के याद रखत रहलें। उनकर भीतर रोशनी के कौंध कबो दिखाई पड़े आ कबो गायब हो जाय। ओह लोग के आपन पिछला जन्म के दर्शन होखत रहलें आ ऊ दिन-रात एह सब के बारे में सोचत-समझत रहलें।\footnote{ज्ञान दीपक, पांडु., पृ. ३९८-३९९} कुछ समय एही तरह बीते पर, ओह लोग के आपन सतगुरु के कृपा प्राप्त भइल आ पवित्र शब्द या शब्द धुन के वरदान मिलल। दरिया एह तथ्य के संदर्भ एह तरह से करे लें:

\begin{quote}
कछु काल बीते जब ज्ञान उपन्नो,\\
सत साहब की कृपा पहिचानो।\\
अति आनंद शब्द सुहावा,\\
प्रेम भगति सुमिरन चलावा।\\
प्रेम सनेह जब प्रगट भयो,\\
गुरु ज्ञान पूरब नाम लयो।\footnote{\textit{दरिया सागर}, चौपाई २६-२७, द.ग्रं. भाग २, पृ. ४}
\end{quote}

बीस बरिस के उमिर में दरिया के भीतरी ज्ञान के पूर्ण प्रकाश भइल। जइसन ऊ कहतारे:

\begin{quote}
बीस वर्ष जब बीति गये,\\
अंतर ज्ञान का घट भरि आयो।\\
सब घट दरसन होइ गयो,\\
जनम मरण ते छूटि गयो।\footnote{\textit{ज्ञान दीपक}, पांडु., पृ. ३९९-४००}
\end{quote}

दरिया आपन परिवार के लोग के विश्वास में लिहलें आ ओह लोग के सतगुरु के माध्यम से गुप्त ज्ञान प्राप्त करे के रास्ता समझवलें। ऊ मांस आ मछली खाए से परहेज आ बिबिध देवी-देवता के पूजा छोड़े के महत्व पर जोर दिहलें। ऊ बतवलें कि संत लोग ही पूजा के योग्य बाड़ें। दरिया बाबा के माई-बाप, भाई आ पत्नी — सभे ओह लोग के उपदेश स्वीकार कइलन।\footnote{वही, पांडु., पृ. ४०४} एही समय ऊ आपन पहिली रचना \textit{दरिया सागर} पूरा कइलन।\footnote{पहिले हम \textit{दरिया सागर} के रचना कइलीं, जेहसे सच्चा योग, वैराग्य आ ज्ञान प्राप्त होला। वही, पांडु., पृ. ४०५}

दरिया बाबा के आपन गाँव आ पड़ोसी गाँव के रूढ़िवादी हिंदू समाज से कड़ा विरोध के सामना करे के पड़ल। पहिला मुकाबला ओह लोग के आपन गाँव के एगो प्रतिष्ठित वैदिक पंडित गणेश उपाध्याय से भइल। दरिया के पहिले ओह लोग से अच्छा संबंध रहे, बाकिर मतभेद तब पैदा भइल जब दरिया मूर्ति-पूजा के निंदा कइलन, आ खासतौर पर पासे के एगो देवी के जानवर बलि के — जेकर पूजा गणेश पंडित करत रहलन आ ओकरा परम देवी मानत रहलन। एह घटना के सजीव वर्णन दरिया के रचना \textit{मूर्ति उखाड़} में मिले ला। दरिया आ गणेश के बीच भइल शुरुआती बातचीत के संक्षिप्त उल्लेख — देवी के परम सत्ता पर गणेश के जोर के संबंध में — बतावे ला कि कइसे ई मतभेद गंभीर रूप ले लिहलस:

\begin{quote}
तब हम (दरिया) जवाब दिहलीं कि सत पुरुष \\
ही परम स्रोत हवें। \\
ई (देवी) त उनकर नौकरानी मात्र बाड़ी। \\
पत्थर काट के मूर्ती बनावल जाला, \\
मूर्ति-पूजा से केहू मूल स्रोत ना पावल। \\
तब ऊ (गणेश) कहलन, "एहन बात कबो ना कहीं। \\
ओकर गोड़ में गिर के पूरा समर्पण करीं।"
तब हम जवाब दिहलीं, "ई मूर्ती त खाली पत्थर के टुकड़ा हवे, \\
चाहीं त तू एकरा के तोड़ सके ला। \\
हाथ, गोड़ आ मुँह सब अंग बनावल गइल बाड़न, \\
बाकिर बोले ना, त सब बेकार बा।"
\end{quote}

गणेश कहलन:

\begin{quote}
"एहन बात तोहरा के कोढ़ी बना देई, \\
आ तोहर दुनों आँखी के नाश कर देई। \\
अब एहिजा से निकल जा आ एक शब्द भी ना बोल \\
अगर गाँव में रहे के चाहत बा।"\footnote{\textit{मूर्ति उखाड़}, पांडु., पृ. २-३}
\end{quote}

गणेश पंडित अब दिन-रात लोगन के दरिया के खिलाफ भड़कावे में लागल रहलन, ओह लोग के समाज खातिर बेहद खतरनाक आ विनाशकारी बतावे लगलन। दरिया के नास्तिक या शैतान करार दिहल गइल जे पहिले कवनो भक्ति में लागल रहलन आ बाद में ओही से ओह लोग के पागलपन हो गइल।\footnote{\textit{ज्ञान दीपक}, पांडु., पृ. ४०६} एह से गणेश पंडित लोगन से गुहार लगवलन कि ऊ लोग आपन भलाई खातिर दरिया के गाँव से निकाल देवे।\footnote{कान दे के हमार बात सुन, आपन हित मन में रख के। ऊ गाँव में ना रहे देई, बाद में तकलीफ पैदा करी। ऊ ना हरि के नाम लेला ना हर के। \textit{ज्ञान दीपक}, पांडु., पृ. ४०७}

एही बीच, दरिया के अत्यंत प्रेमी एगो ब्राह्मण बीरबल जे आपन गाँव के रहनिहार रहलन, एक रात देवी के ओकर जगह से हटा के दूर कहीं जमीन में गाड़ दिहलन। पूरा इलाका देवी के गायब होखे पर सन्न आ व्यथित भ गइल। लोग तरह-तरह के अटकल लगावे लागल। कुछ कहलन कि ऊ कवनो अछूत के छूअ जाए से नाराज हो के चली गइली, चाहे स्वर्ग में या कवनो पाताल में। दूसर लोग कहलन कि ओकरा के कवनो असुविधा या अनादर भइल जेकरा से ऊ जगह छोड़ के चली गइली। अउरी कुछ लोग अटकल लगवलन कि ऊ गुस्सा में बिंध्याचल में आपन बहिन लगे चली गइली होखी।

तीन महीना तक केहू के कवनो सुराग ना मिलल। बाकिर एक दिन दरिया के एगो क्षत्रिय भक्त बाली, जे एह बात से वाकिफ रहे, ताना मार के मामला उगल दिहलन कि बेजान पत्थर के मूर्ती आपन आप ना स्वर्ग जा सके ला ना पाताल, जबले केहू ओकरा के हटावे। एह बात से गाँव वाला लोग समझ गइल कि ओकरा के कुछ पता बा आ ऊ लोग जबरदस्ती कर के ओकरे से सच उगलवलन। अंत में ऊ कहलन कि जा के आपन ताकत असली शक्तिशाली संत दरिया बाबा पर आजमाव।

तुरंत गाँव वालन के एक झुंड दरिया बाबा के चारों ओर से घेर लिहलस आ जोर दे के कहलस कि अगर जान बचावे के चाहत बाड़न त देवी के देखावे के होई। दरिया बाबा निर्भय हो के एह बात से इनकार कइलन जबले कि ऊ लोग गाँव के मुखिया के सामने लिखित प्रतिज्ञा ना दे देवे कि देवी के जानवर बलि ना करे। जब ई शर्त पूरा हो गइल, त दरिया बाबा ओह लोग के ओह जगह देखवलन जहाँ से ऊ लोग आपन देवी के खोद के निकाललस।

बाकिर मामला एहिजा ना रुकल। देवी के खोद के निकारे पर ऊ लोग देखलन कि ओकर नाक टूटल बा। गुस्सल गाँव वाला लोग फैसला कइलस कि देवी के प्रायश्चित करे के एकमात्र उपाय — जे, ओह लोग के अनुसार, तीन महीना से भूखल रहली आ ऊपर से ओकर नाक टूटल — दरिया के ओकरा बलि चढ़ा देवल जाय। एह भयानक काम खातिर, भीड़ दरिया के जगह पर धावा बोल दिहलस आ बलपूर्वक ओह लोग के देवी लगे ले जाय के कोसिस कइलस। बाकिर अचानक दरिया के भाई, भतीजा आ भक्त लोग — जइसे तेग बहादुर, भीखम खाँ, दूँद खाँ, तैयब, दलान, अजीज, चंदन आ अउरी लोग — तलवार आ भाला लहरावत मौके पर पहुँच के गाँव वालन के डरा दिहलस आ ऊ लोग तुरंत भाग खड़ा भइल। दरिया के माई-बाप बहुत डर गइलन, बाकिर दरिया ओह लोग के एह शब्दन में हिम्मत आ साहस बँधवलन:

\begin{quote}
कवनो डर से ना घबरावीं। \\
नाम के दृढ़ता से पकड़ीं \\
जे सदा मदद खातिर तइयार बा। \\
एह दुनिया में आ के \\
हम आत्मा सभ के उद्धार करब \\
शब्द धुन के चीन्हा दे के। \\
ऊ जीवंत प्रभु जे अभिमान तोड़े लें \\
आ पाप मिटावे लें, सदा मदद खातिर मौजूद बाड़न।\footnote{\textit{मूर्ति उखाड़}, पांडु., पृ. १६}
\end{quote}

गाँव वाला लोग, बाकिर, दरिया बाबा के दंड देबे पर अड़ल रहलन। दरिया के सबक सिखावे के पक्का इरादा से, पूरा गाँव ओह शाम मुखिया के जगह पर जमा भइल आ मुखिया दरिया के बोलावे खातिर दूत भेजलस। दरिया तुरंत दूत के संग चल के सभा में आ गइलन। मुखिया गुस्सा में दरिया से मूर्ति तोड़े के सफाई माँगलस आ ओह लोग के शरीर तोड़े के धमकी दिहलस। दरिया निर्भय हो के बेजान मूर्ति के पूर्ण तुच्छता बतवलन। एह से मुखिया के गुस्सा और बढ़ गइल। ऊ आपन तलवार निकाललस, दरिया पर चिल्लावे लागलस आ मारे के धमकी दिहलस। दरिया जान-बूझ के मुखिया लगे बढ़े लागलन। बाकिर जइसे ऊ उग्र मुखिया लगे बढ़लन — जे आपन तलवार उठा के दरिया के मारे खातिर तइयार रहे — अचानक सब लोग के शेर के भयानक गरज सुनाई पड़ल जेकरा से सब हिल गइलन आ डर के मारे मुखिया समेत सब इधर-उधर भाग गइलन।\footnote{वही, पांडु., पृ. १८} दरिया बाबा फिर शांति से घर लौट आइलन।

कुछ समय आपन जगह पर बिता के, दरिया बाबा एक दिन गंगा नदी के तीर पर बसल गाँव बहादुरपुर गइलन। निहाल सिंह जे ओहिजा मौजूद रहलन, दरिया के नदी के तीर पर ठहराव के व्यवस्था कइलस ओह लोग के पवित्र पुरुष मान के। जब दरिया बाबा आपन सामान्य दृढ़ मुद्रा में बइठल रहलन, गणेश पंडित ऊहाँ से हो के गुजरलन। दरिया के पवित्र मुद्रा में बइठल देख के जेकर हाथ में माला, घंटी, धूप या कवनो अउरी पूजा के सामान ना रहे, गणेश पंडित जोर से हँसलन आ कहलन:

\begin{quote}
एह दृढ़ मुद्रा में बइठे से का मिले ला? \\
हम तोहर हाथ में कुछ नइखे देखत। \\
अगर तू गंगा के एहिजा ले आ सके ला \\
तबही हम जानब \\
कि तोहर अभ्यास में कुछ सचाई बा।\footnote{\textit{मूर्ति उखाड़}, पांडु., पृ. २०}
\end{quote}

दरिया कोमलता से जवाब दिहलन:

\begin{quote}
जमीन आ पानी \\
सब प्रभु के आदेश में बा। \\
जहाँ प्रभु आदेश देवे लें, \\
ऊहाँ ऊ दौड़ के आवे ला।\footnote{वही, पांडु., पृ. २०}
\end{quote}

सब गाँव वालन के सामने दरिया के मजाक बनावे के नियत से, गणेश पंडित गाँव में गइलन आ अफवाह फैला दिहलन कि दरिया नदी के तीर पर एह संकल्प के साथ बइठल बाड़न कि जबले गंगा — ओह लोग के नौकरानी — आ के ओह लोग के पवित्र गोड़ ना धो देई, तबले ऊहाँ से ना हटब। केहू साधु के वेश में (संभवतः दरिया के सतगुरु) दरिया के सामने प्रकट भइलन आ ओह लोग से बइठल रहे के आग्रह कइलन। सुबह होखत-होखत भारी भीड़ जमा हो गइल, आ ओह लोग के आश्चर्य के ठिकाना ना रहल जब ऊ लोग गंगा से एगो ऊँच लहर उठत आ शान से दरिया के बइठल जगह के ओर बढ़त देखलस। लोग जोर से हल्ला मचवलस जेकरा से बाकी गाँव वाला लोग, मुखिया समेत, ऊहाँ दौड़ के आ गइलन। तबले लहर वापस गंगा में समा चुकल रहे, बाकिर दरिया बाबा के आसन आ गोड़ गीला रहलन आ ओह लोग के आसपास के जगह पानी से भीगल रहे। सब लोग दरिया बाबा के सामने झुक गइलन आ आपन पहिले के मूर्खता खातिर माफी माँगे लागलन। ई लोगन के दरिया के प्रति नजरिया में निर्णायक मोड़ रहे, आ मुखिया आ गणेश पंडित दुनों दरिया के प्रशंसक बन गइलन।

दरिया के प्रति आपन गाँव में बिरोध, बाकिर, अबहियो बना रहल। आपन आध्यात्मिक अभ्यास में अनावश्यक बाधा से बचे खातिर, दरिया एक दिन आपन प्रिय ब्राह्मण भक्त बीरबल के संग धरकंधा छोड़ के उत्तर के ओर बढ़लन। गंगा नाव से पार कर के, ऊ बलिया जिला के हरदी गाँव पहुँचलन जहाँ गाँव के मुखिया आ ओकर ब्राह्मण पुरोहित ओह लोग के गर्मजोशी से स्वागत कइलन। बाद में, ऊ लोग दरिया बाबा से अनुरोध कइलन कि आपन गाँव वालन के अज्ञानी काम के अनदेखा कर दीं आ सलाह दिहलन कि गाँव लौट जाएं आ ओह लोग पर दया देखावें।

दरिया के चले जाए के बाद, ओह लोग के गाँव वाला लोग भी उदास रहलन आ कुछ पछतावा में रहलन। ऊ लोग सोचे लागलन कि का दरिया सच में एगो पवित्र पुरुष रहलन जिनकर काम ओह लोग के समझ से परे रहे।\footnote{का ऊ (दरिया) एगो सिद्ध ऋषि या संत रहलन? उनकर अद्भुत काम ओह लोग के समझ में ना आइलन। \textit{ज्ञान दीपक}, पांडु., पृ. ४१७}

कुछ समय दरिया अलग-अलग जगह भ्रमण कइलन, आ फिर गंगा पार कर के दूसर किनारे लौट आइलन, बाकिर गाँव से दूरही रहलन। बाद में, ऊ मगहर गइलन आ अयोध्या जाए के सोचलन, बाकिर चूँकि ओह लोग के पत्नी आ माई-बाप ओह लोग के अनुपस्थिति में बहुत उदास रहलन, ऊ लोग लगभग पाँच महीना के यात्रा के बाद ओह लोग के घर वापस भेज दिहलन। दरिया आपन गाँव में लगभग आधा महीना बितवलन।

फिर, क्वार महीना में आसमान में चाँद के उजियारा में, दरिया के सतगुरु (जिनका के दरिया हमेशा सत पुरुष कह के संबोधित करे लें) ओह लोग के जगह पर कृपा कइलन। दरिया के खुशी के ठिकाना ना रहल। ओह लोग के पत्नी पानी से भर के एगो बर्तन ले आइली आ बड़ी भक्ति से सतगुरु के गोड़ धोइलन।\footnote{हमार घरनी घर से बर्तन ले आइली आ पानी से भरलस। ऊ ओह लोग के गोड़ धोइलन आ दर्शन कइलन। \textit{ज्ञान दीपक}, पांडु., पृ. ४१९} दरिया आ ओह लोग के पत्नी के भक्ति से अत्यंत प्रसन्न हो के, सतगुरु कुछ समय ओह लोग के संग रहलन आ दरिया के आध्यात्मिकता के सगरे रहस्य समझवलन। परम आध्यात्मिक धाम के बारे में बतावत ऊ कहलन:

\begin{quote}
परम धाम अकह\footnote{"अकह" शब्द के प्रयोग कबीर भी अनामी पुरुष के परम धाम खातिर कइले बाड़न। — \textit{संतन की बानी} (बियास: राधा स्वामी सत्संग, १९६९), पृ. २२९} (अकथनीय) हवे \\
जे पूरा तरह से भीतरी दर्पण में समायल बा।\footnote{\textit{भक्ति हेतु}, चौपाई २६३, द.ग्रं. भाग २, पृ. ३०३} \\
अकह सच्चा नाम (शब्द धुन) के स्रोत हवे \\
ई (नाम) सबसे सच्चा आध्यात्मिक साधना हवे। \\
एह के जान के, आत्मा आपन चमक पावे \\
आ यम (मृत्यु के राजा) के अभिमान तोड़ दे।\footnote{\textit{भक्ति हेतु}, चौपाई २६३, द.ग्रं. भाग २, पृ. ३०३}
\end{quote}

सतगुरु, बाकिर, सावधानी के एक बात जोड़लन:

\begin{quote}
हम भीतरी रहस्य पूरा तरह से समझा दिहलीं। \\
बाकिर एकर गहिराई के मौन रह के \\
भीतर ही पचावे के होई। \\
एकरा के कहीं प्रकट ना करीं। \\
नाहीं त काल (नकारात्मक शक्ति) के सामने आ जाई।\footnote{\textit{ज्ञान दीपक}, पांडु., पृ. ४२८}
\end{quote}

एक बेर जब सतगुरु अबहियो दरिया के जगह पर रहलन, पूरा इलाका भयंकर सूखा से पीड़ित रहे आ फसल मर रहल रहे। दरिया मन ही मन सतगुरु से अच्छा बरखा खातिर प्रार्थना करे के सोचलन, बाकिर ऊ कुछ कहें से पहिले, सब जाने वाला सतगुरु कोमलता से कहलन कि फसल आ इलाका के लोग खातिर बहुत गरमी हो रहल बा, आ अच्छा बरखा होखे त बढ़िया रही। कुछ देर में भारी बरिश भइल आ लोग बहुत राहत महसूस कइलन। दरिया के बहुत खुशी भइल ई देख के कि ओह लोग के दयालु आ सब जाने वाला सतगुरु ओह लोग के इच्छा पूरा कइलन बिना ओकरा जतावे के।

सतगुरु बाद में दरिया के समझवलन कि सतगुरु कइसे प्यार से आपन शिष्यन के देखभाल करे लें। ऊ दरिया के बतवलन कि जब दरिया खाली एक महीना के रहलन, त ऊ ओह लोग से मिले आइल रहलन, कइसे ऊ पिछला तीस बरिस लगातार दरिया के देखभाल करत रहलन, आ कइसे ऊ मगहर में दरिया के सामने प्रकट भइल रहलन बिना पहिचानल गइल।\footnote{हम तोहरा आगे मगहर में प्रकट भइलीं, बाकिर पहिचान ना करे के कारण तू हमरा के स्रष्टा ना पहिचानलू। जब तू एह दुनिया में जन्मलू, तब हम तोहर जगह आइलीं। तोहर माई तोहरा के हमरा लगे ले आइली आ हम तोहर नाम धरलीं। हम कहलीं कि तोहरा के दरिया कहल जाय आ फिर हम चल गइलीं। हम तोहरा के पिछला तीस बरिस से देखत आइल बानी, आ जहाँ-जहाँ तू गइलू, हम तोहर संग चललीं। \textit{ज्ञान दीपक}, पांडु., पृ. ४४२}

हमनी के आजो ना पता कि ई सतगुरु केन रहलन। दरिया के छोट समय में कई बेर ओह लोग के दर्शन के कृपा मिलल आ अंत में ओह लोग के पूरी गंभीरता के साथ गद्दी (मास्टरशिप) सोंपल गइल। एह घटना के वर्णन \textit{ज्ञान दीपक} में एह तरह से मिले ला:

\begin{quote}
तब सतगुरु आपन घोषणा कइलन, \\
आ एगो सफेद कपड़ा मँगवलन। \\
कपड़ा तुरंत ले आइल गइल, \\
आ सतगुरु के सामने साफ-सुथरा रखल गइल। \\
सतगुरु के निर्देश पर \\
ओकरा सिंहासन पर बिछावल गइल, \\
आ ऊपरी आवरण पूरी गंभीरता से रखल गइल। \\
सतगुरु तब दरिया से एहन कहलन: \\
"कृपा कर के सिंहासन पर बइठीं, \\
ई हम तोहरा के देले बानी।"\footnote{वही, पांडु., पृ. ४४४} \\
दरिया झुक के सिंहासन पर बइठलन। \\
सतगुरु कहलन: "बेटा, अब तोहरा के गद्दी मिल गइल बा।"\footnote{वही, पांडु., पृ. ४४५}
\end{quote}

दरिया के निडर हो के आपन मिशन में आगे बढ़े के उत्साहित करत, सतगुरु कहलन:

\begin{quote}
तोहरा कवनो डर ना होखे के चाही, \\
हम तोहरा के शक्ति दे के ई परंपरा स्थापित कइल बानी। \\
जे लोग अधिकार रखे लें \\
आ जिनका राजा आ रईस कहल जाला, \\
ऊ लोग भी अनुशासित हो जाई \\
अगर तोहर आदेश में बाधा डाले के हिम्मत करन।\footnote{\textit{ज्ञान दीपक}, पांडु., पृ. ४४६}
\end{quote}

दरिया से बिदा लिहले से पहिले, सतगुरु अत्यंत प्रेम आ उदारता से दरिया से कवनो वरदान माँगे के कहलन। ऊ दरिया पर धन आ ऐश्वर्य के वर्षा करे के पेशकश कइलन जे, सतगुरु के अनुसार, खाली संत ही गरिमा आ वैराग्य के साथ संभाल सके ला। बाकिर दरिया के धन या संसारिक ऐश्वर्य के कवनो इच्छा ना रहे आ ऊ सतगुरु के प्रेम आ भक्ति के अलावा कुछ ना चाहत रहलन।\footnote{\textit{मूर्ति उखाड़}, पांडु., पृ. २४-२५; \textit{भक्ति हेतु}, द.ग्रं. भाग २, पृ. ३०८-३०९} बाकिर सतगुरु के आग्रह पर, दरिया ओह लोग से विनती कइलन कि ऊ कबो दूसरन के सामने भीख ना माँगे, आ जे कोई ओह लोग के शिष्य स्वीकार करे, ओकरा के जरूरी भोजन आ कपड़ा मिलत रहे। जइसन दरिया कहतारे:

\begin{quote}
हे दया के सागर, हम तोहर दास हईं; \\
हमार एक बेर के माँग सुन लीं। \\
हमरा के दूसरन आगे हाथ पसार के भीख ना माँगे देईं। \\
जो कुछ तू भेजे ला, \\
ओही से हम तोहर गुणगान करब। \\
बाकिर कृपा कर के भोजन आ कपड़ा के प्रबंध करीं \\
ओह लोग खातिर जिनका तू एह संघ में बोलावे ला।\footnote{\textit{भक्ति हेतु}, चौपाई ३२२, ३२५ आ ३२६, द.ग्रं. भाग २, पृ. ३०९-३१०}
\end{quote}

ई माँग तुरंत स्वीकार करत, सतगुरु कहलन:

\begin{quote}
हम भोजन आ कपड़ा के प्रबंध करब \\
ओह सब खातिर जे संघ में आवे लें। \\
जे कोई तोहरा के उद्धारक मान के \\
तोहरा के पीछे चले, \\
हम ओकरा नरक से बचाइब। \\
हम ओकरा परम धाम में निवास दिहब \\
आ ओकरा दिव्य पुष्प शय्या के \\
आनंद भोगे दिहब।\footnote{वही, चौपाई ३३१-३३३, द.ग्रं. भाग २, पृ. ३१०}
\end{quote}

एह स्पष्ट आश्वासन के साथ, ऊ दरिया के आपन मिशन पूरा करे के सलाह दिहलन:

\begin{quote}
यम तोहर लगे आवे के हिम्मत ना करी, \\
आपन नियत आत्मा सभ के संग परम धाम में उड़ि जा।\footnote{\textit{भक्ति हेतु}, चौपाई २९०, द.ग्रं. भाग २, पृ. ३०६} \\
तोहरा के आत्मा सभ पर मुहर लगावे के अधिकार दिहल गइल बा, \\
आ तू जाने ला कि ई लेन-देन \\
सच्चा नाम (सतनाम) के माध्यम से होला। \\
जे कोई तोहर हाथ के मुहर ले के आवे, \\
हम ओकरा पार लगा दिहब।\footnote{वही, सखी १९, द.ग्रं. भाग २, पृ. ३०६}
\end{quote}

दरिया गृहस्थ आ संन्यासी दुनों के आपन संघ में स्वीकार कइलन, आ खुलेआम घोषणा कइलन कि दुनों के सच्चा नाम के मुहर मिली।\footnote{गृहस्थ आ संन्यासी — दूनो हमार संघ में बाड़न। बाकिर नाम के मुहर बिना केहू कइसे पार जा सके ला? \textit{ज्ञान दीपक}, पांडु., पृ. ५३०} जरूरत एतना बा कि मनुष्य आपन भोजन आ कपड़ा ईमानदार तरीका से जुटा सके। ओही तरह, चाहे केहू सिर ढाँक के रखे या खुला, रास्ता सब खातिर बराबर खुला बा।\footnote{चाहे केहू सिर ढाँक के रखे या खुला, हम दुनों के सच्चा ज्ञान देलें। \textit{ज्ञान दीपक}, पांडु., पृ. ४७५; तुलना करीं वही, पृ. ४५३} साधना भीतरी होखे के कारण, बाहरी रूप-रंग आ ऊपरी रीति-रिवाज के महत्व ना देवे के चाही।

लोगन के डर दूर करे खातिर कि गृहस्थ पूर्ण शुद्धता आ सिद्धता के साथ भक्ति ना कर सके, दरिया खुद ई सवाल उठवलन आ आपन सतगुरु से जवाब पावलन:

\medskip
\noindent{\bfseries दरिया के सवाल:}

\begin{quote}
कृपा कर के ई समझावीं, हे प्रभु, \\
गृहस्थ जीवन में भक्ति कइसे बचल रहे? \\
गृहस्थ काल के जाल से आपन के कइसे बचा सके, \\
आ सबसे ऊँच धाम में जा सके?
\end{quote}

\medskip
\noindent{\bfseries सतगुरु के जवाब:}

\begin{quote}
मामला पर ध्यान से बिचार के, प्रभु कहलन, \\
हमार नाम वास्तव में अनमोल बा।
\end{quote}
एही बात के आगे बढ़ावत दरिया कहतारे:

\begin{quote}
एकाग्रता से ओकरा (नाम) के थाम ले। \\
अइसन मनई लगे काल ना आई। \\
उठत-बइठत ओह पर ध्यान लगावे, \\
आ भीतरी दिव्य ज्योति से प्रेम बढ़ावे। \\
देवी-देवता के सगरी झूठा पूजा छोड़े, \\
आपन असली प्रभु में लीन हो जाए — \\
ओकरा सत्य मान के। \\
उठत-बइठत, परम प्रभु के ध्यान में रहे, \\
आ शब्द धुन में लीन रहे। \\
एह तरह सत्य के शरण ले के, \\
ऊ निश्चय काल के पार हो जाई।\footnote{\textit{ज्ञान दीपक}, पांडु., पृ. ४८६-४८८}
\end{quote}

गद्दी के जिम्मेवारी ले के, दरिया बाबा आपन मिशन पूरी लगन से शुरू कइलन आ ओह लोग के ख्याति आसपास फैलत गइल। पास के जागीर नोखा गढ़ के तत्कालीन सरदार शुजा शाह दरिया बाबा के शिष्य बनलन, आ ओह लोग के अनुरोध पर, दरिया बाबा ओह लोग के राम के असली अर्थ समझवलन — जे प्रचलित राम, दशरथ के पुत्र से अलग बाड़ें, जे दरिया के अनुसार काल या निरंजन के अवतार हवें।\footnote{\textit{ज्ञान रतन}, द.ग्रं. भाग २, पृ. १२५ आ १३५} दरिया बाबा एह विषय पर \textit{ज्ञान रतन} में विस्तार से चर्चा कइले बाड़ें, जे दरिया बाबा आ शुजा शाह के बीच संवाद के रूप में बा।

जइसे-जइसे अधिक से अधिक लोग दरिया बाबा के ओर आकर्षित होखे लगलन, कुछ रूढ़िवादी हिंदू लोग अउरी चिंतित हो गइलन। ओही समय, धर्मदास के एगो वंशज भगवान दास — जिनकर कुछ अनुयायी रहे — धरकंधा आइलन। दरिया बाबा के सतगुरु पहिले से चेतावनी दे दिहल रहलन कि भगवान दास, जे काल के सच्चा एजेंट बाड़ें, के मिले वाला बा आ ओकर मिशन लोग के भ्रमित आ गुमराह करे के बा। लोग भगवान दास के बतवलन कि दरिया के उपदेश प्रचलित हिंदू मान्यता के खिलाफ बाड़े आ कइसे लोग ओह लोग के विकृत उपदेश से गुमराह हो रहल बा। भगवान दास के रिपोर्ट दिहल गइल कि दरिया राम आ कृष्ण के सर्वोच्चता के कमजोर कर रहल बाड़ें आ सतनाम के परम सत्य बता रहल बाड़ें। बतावल गइल कि दरिया, एगो सामान्य मनुष्य होखे के बावजूद, सत पुरुष के अवतार होखे के दावा करे लें आ उपदेश देवे लें कि मुक्ति केवल ओह लोग से सच्चा नाम पा के मिली, ना कि देवी-देवता सभ के पानी-पत्ती चढ़ा के।\footnote{\textit{ज्ञान दीपक}, पांडु., पृ. ५०५-५०७} चूँकि भगवान दास गाँव वाला लोग के प्रचलित विश्वास के समर्थन करत रहलन, एह बात पर सहमति भइल कि पूरा गाँव जमा होखे, दरिया के बोलावल जाय आ ओह लोग से भगवान दास आ गाँव के मुखिया के मौजूदगी में एक बार खातिर आपन स्थिति साफ करे के कहल जाय। कवनो गलत धार्मिक विचार के प्रचार करत पावल गइल पर दरिया के दंड देवे में कवनो बुराई नइखे, ऊ लोग कहलन।

पूरा गाँव जमा भइल आ दरिया के बोलावल गइल। गाँव के मुखिया ओह लोग से सबके सामने बतावे के कहलन कि ऊ काहे माने लें कि राम आ कबीर अलग-अलग बाड़ें, जबकि सब जाने लें कि कबीर आपन पंथ राम के नाम पर स्थापित कइले रहलन।\footnote{वही, पांडु., पृ. ५१५-५१६} दरिया के आपन बात रखे से पहिले, मुखिया भगवान दास से एह बारे में आपन फैसला देवे के कहलन। भगवान दास, पहिले से तय योजना के अनुसार, आपन अंतिम फैसला सुनवलन:

\begin{quote}
राम आ कबीर वास्तव में एक बाड़ें। \\
दू कहल जाला खाली \\
बातचीत के सतही ढंग में।\footnote{वही, पांडु., पृ. ५१६}
\end{quote}

तब दरिया कबीर के रचना सभ के माँग कइलन आ पाठ्य कथन के रोशनी में समझवलन कि भगवान दास के दृष्टिकोण कइसे बहुत सरल बा आ कइसे ऊ कबीर द्वारा स्पष्ट रूप से बतावल राम के सच्चा अर्थ के बिपरीत बा। जब भगवान दास, आपन स्थिति में स्पष्ट विसंगति आ बेतुकापन देखाए के बाद भी, आपन पुरान दृष्टिकोण के दोहरवलन, त दरिया बाबा टिप्पणी कइलन:

\begin{quote}
ऊ एगो ब्यापारी, दुकानदार बा। \\
ऊ डर के मारे सगरी झूठ बोल रहल बा।\footnote{\textit{ज्ञान दीपक}, पांडु., पृ. ५१६}
\end{quote}

गाँव के मुखिया, जे पहिले से दरिया के बिपक्षी रहलन, तुरंत गुस्सा में आगि भइलन आ चिल्ला के कहलन:

\begin{quote}
ऊ डर में बा आ तू नइखस! \\
तू का समझे ला अपना के? \\
हम तोहरा के बाँध के पानी में फेंक देब, \\
आ देखे ला कवन तोहरा बचा सके ला। \\
बार-बार ऊ आपन तलवार पर हाथ रखलस, \\
आ एह तरह आपन भयानक अधिकार देखवलस।\footnote{वही, पांडु., पृ. ५१७-५१८}
\end{quote}

दरिया बाबा निर्भयता से जवाब दिहलन:

\begin{quote}
हमरा लगे खाली दू गो हाथ बाड़े, का तोरा लगे चार बा? \\
का बात जे तू आपन बेकार के घमंड आ अधिकार देखावे ला? \\
अगर तू एतना घमंड से भरल बाड़ू, \\
केहू पेड़ के डाली पकड़ के उखाड़ के देखा।\footnote{वही, पांडु., पृ. ५१८}
\end{quote}

ई तनाव के माहौल में, अचानक एगो अजीब घटना घटल। जइसन दरिया बाबा बतावे लें:

\begin{quote}
सत पुरुष के शक्ति प्रकट भइल, \\
आ भयंकर कोलाहल मच गइल, \\
जइसे कवनो सेना आ गइल होखे। \\
ई कोलाहल बिना कवनो सेना के उठल, \\
केहू एह गुप्त चमत्कार के ना समझ सके।\footnote{वही, पांडु., पृ. ५१९-५२०}
\end{quote}

एह कोलाहल से सब चकित भइलन, आ गाँव के मुखिया संगे सब जने जान बचा के भागलन। एहके बाद, दरिया बाबा के कवनो परेशानी ना दिहल गइल, आ ऊ आठ बरिस तक शांति से धरकंधा में रहलन, जहाँ ऊ ढेर शिष्य बनवलन। आपन मिशन सुचारू रूप से चलावे खातिर, ऊ आपन महत्वपूर्ण शिष्यन के विशेष जिम्मेवारी दिहलन। मुरली दास के सचिव, मणि दास के लेखक, दल दास के रिकॉर्ड कीपर आ वज़ीर दास के निजी सहायक नियुक्त कइल गइल।

ई प्रबंध कर के, दरिया पूरब दिशा में छोटका भ्रमण पर निकललन। रास्ता में, ओह लोग के लहटन गाँव के भिखम दूबे मिललन, जे दरिया बाबा के बड़ा श्रद्धा दिहलन आ ओह लोग के अपना घर ले गइलन। कुछ सवाल-जवाब के बाद जब भिखम दूबे पूरी तरह आश्वस्त हो गइलन कि दरिया सच्चा संत बाड़ें, ऊ आपन चार भाई आ दू गो चचेरा भाई संगे दरिया बाबा के आपन सतगुरु स्वीकार कइलन।

हमनी फ्रांसिस बुकानन के रिपोर्ट से जानित बानी कि मीर कासिम, जे दरिया के समय में बिहार आ बंगाल के नवाब रहलन, दरिया बाबा के १०१ बीघा कर-मुक्त जमीन दान दिहल रहलन। बुकानन शाहाबाद जिला के भ्रमण के दौरान १९१० में धरकंधा गइलन, जब दरिया के उत्तराधिकारी गुना दास के मृत्यु हो चुकल रहे आ ओह लोग के उत्तराधिकारी टेका दास जगह ले चुकल रहलन। बुकानन आपन रिपोर्ट में उल्लेख करे लें कि दरिया संप्रदाय लगे तब १०१ बीघा कर-मुक्त जमीन रहे जे मूल रूप से नवाब मीर कासिम दरिया बाबा के दान दिहले रहलन।\footnote{शाहाबाद रिपोर्ट, पृ. ७८} कहल जाला कि नवाब के भोजपुर अभियान के दौरान, जब ऊ दिनारा में पड़ाव डालले रहलन — जे धरकंधा से कुछ मील के दूरी पर बा — ऊ दरिया बाबा के सम्मान देवे आइलन आ ओह लोग के एगो बहुमूल्य पत्थर उपहार में दिहलन। दरिया सम्मानपूर्वक उपहार लेबे से मना कर दिहलन। बाकिर नवाब जिद कइलन कि ऊ रख लेहें आ बतवलन कि ई कइतना बहुमूल्य बा। नवाब के जाए के बाद, दरिया बाबा ओह पत्थर के पास के तलाब में फेंक दिहलन। बाद में, जब नवाब ई सुनलन, ऊ ओकरा अपमान मानलन आ ओह पत्थर के माँगे दरिया लगे आइलन। कहल जाला कि दरिया बाबा तलाब में हाथ डाल के एकही बेर में कई अइसन पत्थर निकाल लिहलन आ नवाब से कहलन कि आपन पत्थर चुन लेहें। नवाब एह से एतना प्रभावित भइलन कि ऊ दरिया बाबा के १०१ बीघा कर-मुक्त जमीन दान देवे के इच्छा जाहिर कइलन। दरिया बाबा उपहार स्वीकार करे से मना कर दिहलन। बाकिर नवाब जमीन दान देलन आ दानपत्र दरिया के एगो शिष्य के दे के चल गइलन।

दरिया के जीवन के आगे के बिबरण ना मिले ला। बाकिर ई निश्चित बा कि ऊ बहुत लंबा आ अत्यंत सक्रिय जीवन जीहलन, अथक आ निडर हो के आपन मिशन में आगे बढ़त रहलन, भीषण विरोध के बावजूद। बनारस के रामेश्वर पंडित से ओह लोग के लंबा बहस भइल जेकर स्पष्ट उल्लेख उनकर प्रमुख कृति \textit{शब्द} में मिले ला। दरिया कम से कम बीस पुस्तक — संभवतः कुछ अउरी — के लेखक मानल जात बाड़ें। हमनी उनकर औपचारिक शिक्षा के बारे में कवनो जानकारी ना पावित बानी, बाकिर उनकर उत्कृष्ट साहित्यिक शैली, अलंकार के कुशल प्रयोग, सूक्ष्म आ परिष्कृत काव्य उपकरण जइसे कि अनुप्रास, उपमा आ रूपक के प्रभावी उपयोग, आ हिंदी आ उर्दू के बिबिध स्थानीय बोली से लिहल शब्दन के विविधता से ई स्पष्ट बा कि ऊ कवनो भी मानक से भारत के श्रेष्ठ संत-कवियन में से एक रहलें।

लागत बा कि ऊ अपना जीवनकाल में अपना संप्रदाय के कुछ उपकेंद्र स्थापित कइले रहलन। विक्रम संवत १८३६ (१७७९ ई.) में, यानी आपन शरीर छोड़े से कुछ महीना पहिले, ऊ अपना सबसे उन्नत शिष्यन में से दू गो — गुना दास आ टेका दास — के अपना उत्तराधिकारी नियुक्त कइलन। दरिया बाबा ओह लोग के आपन आध्यात्मिक पुत्र कहत रहलन आ ओह लोग के संग बड़ा प्रेम से बराबरी के व्यवहार करत रहलन।\footnote{द.ग्रं. भाग १, पृ. ३; \textit{मूर्ति उखाड़}, पांडु., पृ. २९} दरिया बाबा ई भी स्पष्ट कइलन कि एह उत्तराधिकारी लोग के अपना उत्तराधिकारी नियुक्त करे के अधिकार होई, आ एही तरह ई परंपरा चलत रही।\footnote{\textit{मूर्ति उखाड़}, पांडु., पृ. ३०}

जब ओह लोग के भाई फक्कर दास पूछलन कि ई परंपरा कब तक चली, दरिया बाबा जवाब दिहलन कि जब तक शब्द धुन के अभ्यास संप्रदाय के केंद्र में रही, आ संप्रदाय शुद्ध आ बाहरी रीति-रिवाज आ आडंबर से मुक्त रही, तब तक ई परंपरा चलत रही। जब शब्द धुन के अभ्यास बाहरी रूप (भेख) आ बाहरी रीति-रिवाज में मिल जाई, तब शब्द धुन अलग हो जाई। तब हमार द्वारा दिहल शक्ति चली जाई, आ आत्मा सभ काल के मुँह में चली जाईं। ओह अवस्था में, ऊ एह संसार में आवे आ शब्द धुन के एगो नवा परंपरा शुरू करी। एही तरह, ऊ कहलन, ऊ युगों से आवत रहल बाड़ें।\footnote{वही, पांडु., पृ. ३१-३३}

एह तरह, संत मिशन के सगरी बात आ सगरी पहलू समझा के, आ आपन जीवन के उद्देश्य पूरा कर के, ऊ आपन जिम्मेवारी आपन उत्तराधिकारी लोग के सौंप के विक्रम संवत १८३७ (१७८० ई.) में एह संसार से विदा हो गइलन।\footnote{वही, पांडु., पृ. ३१-३३}

उपरोक्त बिबरण से ई स्पष्ट बा कि दरिया सच्चा संत जीवन के एगो गौरवशाली उदाहरण रहलें, जे पूरी तरह से मानव जाति के सर्वोच्च सेवा में समर्पित रहलें। दया के मिशन के पूरा करे में, ऊ आपन जान के जोखिम में डाले में कबो हिचकिचवलें ना। सबसे भीषण विरोध के बीच भी, ऊ पूरी तरह निर्भय रहलें। एकर परिणाम ई भइल कि ऊ अंधविश्वास आ बाहरी रीति-रिवाज के खोखलापन उजागर करे में सफल रहलें आ आपन दिव्य संदेश बिहार आ उत्तर प्रदेश के बिसाल क्षेत्र में फैला दिहलन। अइसन समय में जब एगो संत खातिर पैदल यात्रा ही परिवहन के एकमात्र साधन रहे, दरिया के उपलब्धि वास्तव में असाधारण रहे।

दरिया जेकरा खातिर खड़ा रहलन, ऊ जे कइलन, उनकर जीवन-शैली आ ओह लोग के विरोधियन के तीखा आ व्यंग्यात्मक जवाब के देख के, ओह लोग के एगो सिद्ध मालिक होखे के दावा वाकई विश्वसनीय लागे ला। जे भी हो, ए बात से इनकार नइखे कि सगरी सच्चा संत मूल रूप से दिव्य शब्द (शब्द धुन) के प्रकटीकरण हवें, आ ठीक एही बात दरिया आपन बारे में दावा करे लें। जइसन ऊ कहतारे:

\begin{quote}
हम एह संसार में आइल बानी \\
दिव्य शब्द (शब्द धुन) के प्रकटीकरण के रूप में, \\
आ मानव रूप धारण कर के, \\
हम प्रभु के महिमा आ गरिमा के गावल बानी।\footnote{वही, पांडु., पृ. ४४}
\end{quote}

% अध्याय १ के अंत
